% Review draft only. Not integrated into the delivered manuscript.
\subsection{\bt{Observed reuse behavior}{观察到的复用行为}}
\label{sec:sharing-eval}
\input{figures/reuse-work}
\bi{\emph{Work distribution.} Figure~\ref{fig:reuse-work} describes agent-issued calls under the complete evaluated policies. Over reuse cutoffs, \system\ makes \OfStackExplore{} schema/probe calls per session on rel-stack versus \OfStackNoneExplore{} without memory, and \OfTrialExplore{} versus \OfTrialNoneExplore{} on rel-trial. The validation bars count the common feature-validation tool, whose prediction scoring and temporal audit apply to every method. They do not count \system's publication gates or isolate the effect of shared condition checks. Internal memory-model work is charged in Table~\ref{tab:official-results}. Thus this figure describes changes in agent behavior, not a component ablation of executable definitions.}
{\emph{工作量分布。}图~\ref{fig:reuse-work}描述完整受评估策略下智能体发起的调用。复用截点上，\system\ 在rel-stack每会话进行\OfStackExplore{}次模式检查或SQL探查，无记忆为\OfStackNoneExplore{}次；rel-trial分别为\OfTrialExplore{}与\OfTrialNoneExplore{}次。验证柱统计共同的特征验证工具，其预测评分与时间审计适用于所有方法；这些调用不代表\system{}的发布门控，也未分离共享条件检查的作用。记忆模块内部的模型工作计入表~\ref{tab:official-results}。因此，本图描述智能体行为变化，并非可执行定义的组件消融。}

\bi{\emph{Tracing the consumed representation.} Among the twelve reuse sessions, nine receive a complete previous task SQL and validate that SQL first; six retain it verbatim as their final SQL. No reuse session directly validates an unchanged definition-compiled SQL. Neither database records an explicit \code{compose\_features} call, but this alone does not imply that compilation is absent: \code{find\_features} and \code{find\_program} also compile definitions internally. Returning a compiled SQL or incrementing a definition's lookup counter does not establish that the agent executes or selects it. Manual modifications may incorporate retrieved knowledge; the current traces do not identify their separate contribution.}
{\emph{追踪实际使用的表示。}十二个复用会话中，九个取得上一任务的完整SQL并首先验证它，六个将其逐字保留为最终SQL。没有复用会话直接验证原样返回的定义编译SQL。两个数据库均无显式\code{compose\_features}调用，但这不代表完全没有编译：\code{find\_features}和\code{find\_program}也在内部编译定义。返回编译SQL或增加定义命中计数，并不能证明智能体执行或选用了该结果。人工式改写可能吸收检索到的知识，现有轨迹不能分离它的独立贡献。}

\bi{\emph{A transfer limitation.} At the first cutoff, outside Figure~\ref{fig:reuse-work}'s reuse subset, user-engagement directly validates eight features compiled from the shared library. The common temporal audit rejects three post-related columns, despite their publication checks having passed on the source user-badge task. The agent subsequently constructs a different SQL that passes the audit. This observation limits the scope of the publication evidence: checks on source instances do not certify every target-task population. The official-API comparison therefore establishes the behavior and cost of a policy combining program reuse, optional definition retrieval, and renewed validation; it does not establish an incremental quality or safety benefit from feature composition. A controlled representation ablation on held-out targets is required for that claim.}
{\emph{迁移的限制。}首截点的user-engagement会话不属于图~\ref{fig:reuse-work}的复用子集，它直接验证共享库编译出的八个特征。共同时间审计拒绝了其中三个与帖子有关的列，尽管它们在来源user-badge任务上的发布检查均已通过。智能体随后构造了另一份通过审计的SQL。这一观察限定了发布证据的适用范围：来源样本上的检查不能认证所有目标任务总体。因此，官方API比较确立的是整份方案复用、按需定义检索和重新验证的组合策略的行为与代价，尚未确立特征组合带来的额外质量或安全收益。该主张需要在留出目标上进行受控的表示消融。}

\bi{\emph{Publication fidelity in the metric studies.} The separate controlled metric workload checks whether extraction preserves verified answers. Across \PubAllJudged\ verified sessions, SQL replay rejects \PubReplayRefused\ erroneous candidates and the calculation check rejects \PubChainRefused\ more. These observations assess publication fidelity independently of the RelBench reuse-cost comparison; shared maintenance and snapshot validity are evaluated separately in Sections~\ref{sec:cost} and~\ref{sec:changes-eval}.}
{\emph{指标实验中的发布保真性。}独立的受控指标负载检验提取是否保留已核实的答案。在\PubAllJudged{}个答案已核实的会话中，SQL重放拒绝\PubReplayRefused{}个错误候选，算式检查另拒绝\PubChainRefused{}个。这些观察独立于RelBench复用成本比较，检验发布保真性；共享维护和快照有效性分别在第~\ref{sec:cost}、\ref{sec:changes-eval}节评估。}
